%=======================================================================================
% DATA CHAPTER -- 2026-08-12, 18:10. Visual-first restructure.
%
% Supersedes documentation/august_draft/3_Chapter_data/data_chapter_merged_2026-08-12_1706.tex
% (tables-as-placeholders version) and data_chapter_merged_2026-08-12.tex (fully-typeset-tables
% version). Neither earlier file is deleted.
%
% What changed this pass, structurally (content/numbers unchanged unless flagged \fc below):
%   - Old S1 (Study area and source data) + old S2 (Spatial and temporal hierarchy) merged into
%     one opening section. The hierarchy detail (block/compartment/sub-compartment correction)
%     moves into the opening figure's Panel C instead of its own prose section.
%   - The old free-standing cohort-comparison table is removed; its content (final plots/records/
%     compartments/role) is folded into the cohort-coverage map's caption instead, per the
%     instruction not to keep a second table when a figure caption already carries the numbers.
%   - The cleaning funnel table is collapsed from 7 rows to 5 (the three plausibility checks --
%     planting year, age, height -- become one table note, since they remove ~0 additional rows
%     on the current data; verified again this pass, unchanged from the previous version).
%   - The itemised list of original stand/management fields becomes a table (Field | Definition |
%     Modelling status); the *reasoning* behind each flagged field (CanopyCover, LAI, yldc) stays
%     in prose, since compression rule 3 keeps "why" in text, not in a table cell.
%   - Strengths/limitations becomes one paired table instead of two prose paragraphs.
%   - The repeated-observation-timeline figure is dropped: the five gap lengths are a short enough
%     sequence to state once in a sentence, not a figure of their own.
% Colour-coded flags: \fc = fact-check (red), \ct = citation missing (orange), \app = move to
% appendix (blue), \sr = small reword (teal) -- same legend as 1_Chapter_Intro. No model results,
% performance numbers, or causal language appear in this chapter.
%=======================================================================================
\newcommand{\fc}[1]{\textcolor{red}{[FactC: #1]}}
\newcommand{\ct}[1]{\textcolor{orange}{[CITE: #1]}}
\newcommand{\app}[1]{\textcolor{blue}{[APP: #1]}}
\newcommand{\sr}[1]{\textcolor{teal}{[Reword: #1]}}
\newcommand{\fig}[1]{\textcolor{purple}{[FIG: #1]}}
%=======================================================================================

\chapter{Data}
\label{ch:data}

This chapter describes the data on which every result in this dissertation is built. It states
what exists and how it was assembled -- study area, plot geometry, the two analytical cohorts, the
response variable, original stand fields, and the external environmental data matched to each plot
-- not how any of it was selected, split or modelled for a specific research question, which is
described in Chapter~\ref{ch:methodology}. Spatial layouts, filtering stages and comparisons are
shown in figures and tables; prose is reserved for why a choice was made and what it limits.

\section{Study area, plot geometry and cohort coverage}
\label{sec:data-source}

The study area is Aberfoyle forest, Scotland, and the plot records span 21 years (2002--2023)
across a study extent of approximately 16\,km by 12\,km (British National Grid, EPSG:27700). Each
row in the source data represents one plot in one LiDAR survey year; the LiDAR-derived fields
include height percentiles and canopy cover, computed by Forest Research's own Area-Based-Analysis
pipeline from Airborne Laser Scanning (ALS) data. Most plots are approximately 20\,m by 20\,m; on
the four-survey cohort, 76.2\% of plot-year rows are exactly 400\,m\textsuperscript{2}, with the
remainder smaller (clipped by a compartment or sub-compartment boundary) or, in a smaller share,
larger than the nominal cell. \fc{A figure of "28.8\% clipped" was supplied for this section but
does not exactly match direct recomputation this pass: 23.8\% of four-survey rows are not exactly
400\,m\textsuperscript{2} (76.2\% are), and 28.9\% at the raw, pre-cleaning layer (795,645 rows).
Close, but not identical -- worth confirming which population (raw layer, four-survey cohort, or
unique plots rather than plot-years) the 28.8\% figure was computed on before it is quoted as the
headline number. The remainder of this section reports the directly recomputed figures.} The
full spatial layout, the plot/compartment/sub-compartment hierarchy, and examples of regular versus
clipped plots are shown in Figure~\ref{fig:data-study-area}, not repeated here.

Exact ALS acquisition metadata -- flight date, point density, sensor -- was not found in any file
inspected for this dissertation. Because survey years are processed independently, any undocumented
difference in acquisition conditions between years is a source of unmeasured variation in height
observations that this dissertation cannot separate from real growth \ct{Forest Research ALS/LiDAR
acquisition specification, if a technical report exists}.

The analysis is restricted to Sitka spruce (\texttt{spis == "SS"}, 56\% of raw plot-year records),
the dominant species at this site and in British commercial forestry generally, and the only
species with enough repeated, balanced observations to support the cohort construction in
\S\ref{sec:data-cohorts}. This restriction is a scope decision, not a claim that Sitka spruce is
immune to the site's principal risk factors: it is comparatively more wind-firm and productive on
exposed upland sites than most alternative commercial conifers grown in Britain
\ct{comparative Sitka spruce wind-firmness/productivity citation, e.g.\ a Forest Research species
comparison}, but remains susceptible to windthrow, root disease and other biotic and abiotic
pressures like any planted species; no disease- or pest-immunity claim is made anywhere in this
dissertation.

% *** FIGURE SPEC STATUS: CURRENT as of 2026-08-13, user-revised -- 3 PANELS ONLY (A/B/C). ***
% *** The clipped-vs-regular plot example that used to be a separate Panel D is now an annotation ***
% *** INSIDE Panel C. Any future edit to this figure must reconcile points 1-8 together -- do not ***
% *** edit one point (e.g. add a panel, change a layer) without updating the others below to match.***
% FIGURE PLACEHOLDER (study area, hierarchy and plot geometry) -- fully specified
% 1. PURPOSE: orient the reader to where Aberfoyle is, how much of it is covered, how the
%    plot/compartment/sub-compartment hierarchy nests, what species make up the plots, and what a
%    "regular" versus "clipped" plot actually looks like on the ground -- replacing four
%    paragraphs of spatial description.
% 2. PANELS (3 panels: A, B, C -- no Panel D, see status note above):
%    Panel A -- Aberfoyle's position within Scotland. A simplified Scotland outline with a small
%      red marker/box at Aberfoyle's location (Stirling council area, the Trossachs). No plot data
%      needed; a single labelled point is sufficient. Purpose: national context only.
%    Panel B -- the full study extent, on a satellite/aerial image base layer (NEW DEPENDENCY --
%      not currently cached anywhere in this project; needs an external tile fetch at
%      figure-build time, e.g. contextily/ESRI World Imagery or an OS Maps aerial layer -- flag
%      for whoever builds this figure, since every other layer in this chapter comes from files
%      already in the repo and this one does not). Bounding box from plot coordinates in
%      data/processed/environmental/plot_environmental_features.parquet (x 239,327-256,111;
%      y 695,230-707,244; EPSG:27700). Layer order, bottom to top: satellite image -> compartment
%      boundaries (data/interim/compartment_boundaries.parquet, light/thin line) -> block
%      boundaries (data/interim/block_boundaries.parquet, the outermost/heaviest outline). Scale
%      bar in km; north arrow.
%    Panel C -- local enlargement (one representative block or block-pair, chosen for a visually
%      clear compartment/sub-compartment mix and a visible clipped plot). Layers, all from
%      data/interim/*_boundaries.parquet: compartment boundaries (solid line, one colour),
%      sub-compartment boundaries (dashed line, a second colour) -- labelled to make explicit that
%      a sub-compartment letter (A-Z) is only unique in combination with its parent compartment,
%      not on its own (3,385 distinct (cpmt,scpt) pairs exist project-wide, vs. only 26 raw scpt
%      letter values) -- and individual plots as small filled squares/points, coloured by species
%      (Sitka spruce vs. other species present in this block, from the raw layer's \texttt{spis}
%      field; the broadleaf/conifer grouping used for "other" is not yet decided -- needs a
%      species-code lookup this project does not currently have codified anywhere). Cohort
%      membership is NOT shown in this panel (already covered by the separate cohort-coverage
%      figure, \S\ref{sec:data-cohorts}) -- do not reintroduce cohort colouring here, it would
%      duplicate that figure. One arrow-annotated callout within this panel: a regular 20 m x 20 m
%      plot cell (400 m^2) next to a cell clipped by a compartment/sub-compartment boundary
%      (250.0 m^2, the observed minimum), each labelled with its exact area -- this replaces the
%      old separate "Panel D." Scale bar in metres (e.g.\ 200 m).
% 3. DATA REQUIRED: data/interim/block_boundaries.parquet, compartment_boundaries.parquet,
%    subcompartment_boundaries.parquet; data/processed/environmental/plot_environmental_features.parquet
%    (x/y); a species field (\texttt{spis}) per plot, from
%    data/processed/master/clean_master_4survey.parquet or the raw GPKG, for Panel C's plot
%    colouring; data/raw/LiDAR_Years_All_7jul.gpkg read WITH geometry, for the two real plot
%    polygons in Panel C's clipped-vs-regular callout (a new read path not used elsewhere in the
%    pipeline, since the master-export pipeline drops geometry via ignore_geometry=True); an
%    external satellite/aerial basemap tile source for Panel B (see Panel B note above -- not
%    presently a project dependency).
% 4. COLOUR/SYMBOL RULES: one consistent colour per hierarchy level across panels B/C (block =
%    grey/heavy outline, compartment = solid dark line, sub-compartment = dashed line, in Panel B
%    compartment lines are drawn lighter/thinner than in Panel C since B shows many at once);
%    Panel C's plots coloured by species (Sitka spruce vs. other, see point 2); the two callout
%    polygons in Panel C use a third, high-contrast outline colour so they read as an inset
%    example, not more of the same plot layer.
% 5. ANNOTATIONS/LABELS: scale bar and north arrow on B and C; a small inset box on A showing
%    where B's extent sits within Scotland; in Panel C, an arrow/callout labelling the regular vs.
%    clipped plot pair directly with their areas in m^2, not only in the caption; a small colour
%    legend in Panel C for the species categories shown.
% 6. CAPTION MUST STATE: EPSG:27700; extent ~16.8 km E-W x 12.0 km N-S; 18 blocks, 531 raw
%    compartments, 3,385 (cpmt,scpt) sub-compartments (the true unique key, not the 26-letter scpt
%    field alone); median plot area exactly 400 m^2; 76.2\% of four-survey rows exactly 400 m^2,
%    remainder clipped smaller (down to 250.0 m^2) or, less commonly, larger; what the species
%    colours in Panel C represent.
% 7. REMOVE FROM MAIN TEXT once this figure exists: the numeric hierarchy counts (18/531/3,385),
%    since they belong in the caption only -- the main text above states only that the hierarchy
%    and geometry are shown in the figure.
% 8. DRAFT CAPTION: "Aberfoyle study area (EPSG:27700, extent 16.8 km x 12.0 km) within Scotland
%    (A), full plot coverage on a satellite base layer (B), and a local enlargement showing the
%    block/compartment/sub-compartment hierarchy and plot locations, coloured by species (C):
%    18 blocks, 531 raw compartments, and 3,385 (compartment, sub-compartment) pairs -- the
%    sub-compartment field alone takes only 26 letter values (A-Z) reused across many
%    compartments, so compartment and sub-compartment together, not sub-compartment alone,
%    identify a unique management unit. The callout in (C) shows two real plot polygons: a
%    regular 20 m x 20 m cell (400 m^2, the median) and a cell clipped by a sub-compartment
%    boundary (250.0 m^2, the observed minimum)."
\begin{figure}[!htbp]
\centering
\fbox{\parbox{0.9\textwidth}{\centering\vspace{5cm}FIGURE PLACEHOLDER --- study area, hierarchy
and plot geometry (4 panels), fully specified in the comment above source\vspace{5cm}}}
\caption{Aberfoyle study area, spatial hierarchy and plot geometry. See source comment for the
complete panel specification.}
\label{fig:data-study-area}
\end{figure}

\app{Full raw column inventory (all 38 fields in \texttt{LiDAR\_Years}), including fields never
used downstream (\texttt{flyr}, constant at zero for every row; \texttt{block}, a
99.999\%-identical duplicate of \texttt{blk}).}

\section{Cohort construction}
\label{sec:data-cohorts}

Two analytical cohorts were built from the same raw layer, differing only in which survey years
are required to be present for every retained plot: a \textbf{four-survey} cohort (2008, 2012,
2021, 2023) and a \textbf{six-survey} cohort (2002, 2006, 2008, 2012, 2021, 2023 -- unevenly
spaced, with gaps of 4, 2, 4, 9 and 2 years). Both cohorts retain only balanced trajectories: a
plot is excluded if any one of its required survey observations is absent.

% TABLE (final, not a placeholder -- exact figures re-verified 2026-08-12 against
% data_processing/clean_master_data.py::clean_cohort, live re-execution)
\begin{table}[!htbp]
\centering
\scriptsize
\setlength{\tabcolsep}{3pt}
\renewcommand{\arraystretch}{0.95}
\begin{tabular}{p{4.2cm}p{2.4cm}p{2.4cm}p{2.4cm}p{2.4cm}}
\hline
Stage & 4-survey rows & 4-survey plots & 6-survey rows & 6-survey plots \\
\hline
Raw dataset & 795,645 & 230,359 & 795,645 & 230,359 \\
Selected survey years & 735,219 & 230,345 & 795,645 & 230,359 \\
Complete survey coverage & 451,128 & 112,782 & 112,686 & 18,781 \\
Sitka spruce & 289,580 & 72,395 & 83,382 & 13,897 \\
Final valid cohort & 287,064 & 71,766 & 83,382 & 13,897 \\
\hline
\end{tabular}
\caption{Cleaning funnel from the raw layer to the two final cohorts. \emph{Final valid cohort}
additionally requires: a planting year no earlier than 1850 and before the survey year; age
between 0 and 200 years; height greater than 0\,m and at most 60\,m. These three checks are wide
sanity bounds, not domain-derived thresholds, and remove almost nothing on the current data (0
additional rows in the six-survey cohort; a small number in the four-survey cohort, absorbed into
the totals above) -- nearly all attrition comes from requiring complete survey coverage and
restricting to Sitka spruce.}
\label{tab:data-funnel}
\end{table}

% FIGURE PLACEHOLDER (cohort spatial coverage + cleaning funnel, combined side-by-side) -- fully specified
% 1. PURPOSE: combine the cohort coverage map and the cleaning-funnel table into one side-by-side
%    figure to save page space, since both describe the same two cohorts.
% 2. LAYOUT: two-column layout inside one figure float -- left ~56% width: the coverage map;
%    right ~40% width: the cleaning-funnel table, set in \scriptsize so it fits alongside the map.
% 3. LEFT PANEL (map): Aberfoyle/block outline (data/interim/block_boundaries.parquet) as the
%    boundary; sub-compartment boundaries (data/interim/subcompartment_boundaries.parquet) drawn
%    and coloured individually (one categorical colour per sub-compartment, or a light categorical
%    fill if the full palette is too busy across ~150+ sub-compartments) so coverage is visible at
%    that level; four-survey plots as one marker/colour, six-survey plots overlaid as a second,
%    darker marker/colour on top (six-survey is a complete subset, so the overlay shows exactly
%    which part of the coverage it retains). Scale bar, north arrow.
% 4. RIGHT PANEL (table): Table~\ref{tab:data-funnel}'s five stages reproduced in \scriptsize --
%    Stage / 4-survey rows / 4-survey plots / 6-survey rows / 6-survey plots -- so the map and the
%    funnel that produced its two cohorts sit on the same page.
% 5. COLOUR/SYMBOL RULES: one categorical colour per sub-compartment on the map, falling back to
%    cohort-only shading if a full sub-compartment palette proves illegible; four-survey vs.
%    six-survey plot markers as previously specified.
% 6. CAPTION MUST STATE: final plots -- 71,766 (4-survey) vs. 13,897 (6-survey); final plot-year
%    records -- 287,064 vs. 83,382; observations per plot -- 4 vs. 6; compartments represented --
%    296 vs. 48; analytical role -- main analysis vs. sensitivity analysis; the nesting fact --
%    every one of the 13,897 six-survey plots is also a four-survey plot (0 six-survey-only
%    plots), a complete, not partial, subset.
% 7. REMOVE FROM MAIN TEXT once this figure exists: the standalone cohort-comparison table:
%    figures already replaced in this draft (see banner comment at top of file).
% 8. DRAFT CAPTION: "Left: spatial coverage of the four-survey cohort (71,766 plots, 287,064
%    records, 296 compartments, main analysis) and the six-survey cohort (13,897 plots, 83,382
%    records, 48 compartments, sensitivity analysis), sub-compartments coloured individually
%    within the Aberfoyle boundary. Every six-survey plot is also a four-survey plot -- a
%    complete, not partial, subset. Right: the cleaning funnel producing both cohorts from the
%    same 795,645-record raw layer (full validity-check detail: Table~\ref{tab:data-funnel})."
\begin{figure}[!htbp]
\centering
\begin{minipage}{0.56\textwidth}
\centering
\fbox{\parbox{0.95\textwidth}{\centering\vspace{4cm}MAP PLACEHOLDER --- Aberfoyle outline,
sub-compartments coloured individually, 4-survey vs.\ 6-survey plot coverage overlaid, see
comment above source for full specification\vspace{4cm}}}
\end{minipage}%
\hfill
\begin{minipage}{0.40\textwidth}
\centering
\scriptsize
\setlength{\tabcolsep}{2pt}
\renewcommand{\arraystretch}{0.9}
\begin{tabular}{p{2.1cm}p{0.85cm}p{0.85cm}p{0.85cm}p{0.85cm}}
\hline
Stage & 4s rows & 4s plots & 6s rows & 6s plots \\
\hline
Raw dataset & 795,645 & 230,359 & 795,645 & 230,359 \\
Selected years & 735,219 & 230,345 & 795,645 & 230,359 \\
Complete coverage & 451,128 & 112,782 & 112,686 & 18,781 \\
Sitka spruce & 289,580 & 72,395 & 83,382 & 13,897 \\
Final valid cohort & 287,064 & 71,766 & 83,382 & 13,897 \\
\hline
\end{tabular}
\end{minipage}
\caption{Left: four-survey and six-survey cohort spatial coverage (sub-compartments coloured
individually within the Aberfoyle boundary; 71,766 vs.\ 13,897 plots, 287,064 vs.\ 83,382 records,
296 vs.\ 48 compartments, main analysis vs.\ sensitivity analysis; the six-survey cohort is a
complete, not partial, subset). Right: the cleaning funnel producing both cohorts from the same
795,645-record raw layer (full validity-check detail: Table~\ref{tab:data-funnel}).}
\label{fig:data-cohort-coverage}
\end{figure}

The six-survey cohort's completeness as a subset, not merely its overlap, is what makes it a
genuine sensitivity check on the same underlying plots rather than an independent replication --
how that check is actually used for a given research question is described in
Chapter~\ref{ch:methodology}.

\section{Response variable and original stand fields}
\label{sec:data-response}

The response variable used throughout this dissertation is \texttt{elev\_percentile\_95th}: the
95th percentile of ALS return heights above ground within a plot, for a given survey year, kept
in its raw form with no correction applied. Table~\ref{tab:data-fields} lists this and every other
original stand and management field carried into the master export, together with its modelling
status; the reasoning behind the three flagged rows follows in prose below the table.

% TABLE (final, not a placeholder)
\begin{table}[!htbp]
\centering
\scriptsize
\setlength{\tabcolsep}{3pt}
\renewcommand{\arraystretch}{0.95}
\begin{tabular}{p{2.6cm}p{6.3cm}p{3.9cm}}
\hline
Field & Definition / derivation & Modelling status \\
\hline
\texttt{elev\_percentile\_95th} & 95th percentile of ALS return heights above ground; raw, no correction & Target (every RQ) \\
\texttt{Top\_Height95} & $=\texttt{elev\_percentile\_95th}\times1.1$ (exact, verified) & Audit ingredient only (feeds \texttt{Vol95}/\texttt{GYCspec95}); never a feature \\
\texttt{Top\_Height99}, \texttt{Vol99}, \texttt{GYCspec99}, \texttt{elev\_percentile\_99th} & 99th-percentile family & Retired 2026-07-28; not used anywhere \\
\texttt{plyr}, \texttt{Age} & Planting year; \texttt{Age} = survey year $-$ \texttt{plyr} & Modelling input; also a cleaning-stage filter \\
\texttt{yldc} & FR yield class, 12 distinct values (2--24) & Audit/stratification only -- flagged$^{\dagger}$ \\
\texttt{whcl} & Ordinal FR windthrow hazard class, 0--6 (6 of 7 levels observed) & Management assessment, not a wind measurement \\
\texttt{CanopyCover} & ALS-derived canopy-cover fraction (0--1) & Baseline environmental feature -- flagged$^{\ddagger}$ \\
\texttt{Thin}, \texttt{last\_thinn}, \texttt{next\_thin\_} & Thinning history; derives \texttt{time\_since\_thinning} and recency flags & Modelling input (\app{exact derivation formulas}) \\
\texttt{LAI} & ALS-derived leaf area index & Kept, unused as a feature -- flagged$^{\S}$ \\
\texttt{Vol95}, \texttt{GYCspec95}, \texttt{area}, \texttt{p1--p5}, \texttt{g1}, \texttt{g3} & FR reporting/audit fields; several are direct functions of height and age \ct{Forest Research, "Use of Airborne Laser Scanning (ALS) for Forest Inventory" -- source of p1--p5, full reference not yet confirmed} & Reporting only; never predictors (\app{exact formulas}) \\
\hline
\end{tabular}
\caption{Response variable and original stand/management fields kept in the master export.
$^{\dagger\ddagger\S}$ reasoning below.}
\label{tab:data-fields}
\end{table}

$^{\dagger}$\textbf{Yield class} is excluded from every candidate feature set because it is itself
derived from forest growth information: including it alongside a height-based target or a
growth-curve-residual target raises a leakage/circularity concern.

$^{\ddagger}$\textbf{Canopy cover} is a live, checked question rather than a settled exclusion.
\fc{\texttt{CanopyCover} and the height target come from the same GeoPackage row -- the same ALS
flight and area-based-analysis pipeline, for the same plot and survey year -- unlike the
terrain/wind/climate/soil variables in \S\ref{sec:data-environmental}, which come from entirely
independent sources. This is a reverse-causation question, not an algebraic leak like
\texttt{yldc}'s: canopy closure could be a consequence of prior growth, a partial independent
driver of it (via planting density, stocking or thinning regime), or both, and shared provenance
alone does not distinguish these. \texttt{CanopyCover} is nonetheless included as one of four
baseline features in every environmental feature set used across all three research questions
(\texttt{documentation/methodlogy\_env\_setpick.md}, \S1.6/\S2.11). This is addressed empirically,
not by assumption: a baseline-only ablation
(\texttt{TEMP\_results/TEMP\_rq2\_attribution\_results\_2026-08-11.tex}, "CanopyCover
reverse-causation check," added 2026-08-13) fits RQ2's attribution models on the baseline columns
alone and compares against the full environmental sets. This chapter states the shared-provenance
fact only; the ablation's outcome and its consequences for how \texttt{CanopyCover} is framed
belong in Chapter~\ref{ch:methodology}'s feature-set justification and the Results chapter --
verify those sections carry this through once drafted.}

$^{\S}$\textbf{LAI} shows a clear processing discontinuity in 2021 specifically -- median 3.09 and
maximum 23.0 that year, against a median of 0.4--0.9 and a hard ceiling of 5.5 in every other
survey year (verified directly on the raw layer, grouped by survey year) -- consistent with a
different LAI-processing algorithm for the 2021 acquisition. Since LAI is not used as a model
input anywhere in this project, this does not affect any result in this dissertation, but it is
stated here so the column is not silently treated as trustworthy if reused later.

% FIGURE PLACEHOLDER (target and age coverage) -- fully specified
% 1. PURPOSE: show the range and density of the underlying age/height data across cohorts and
%    survey years, without presenting any fitted model or result.
% 2. PANELS: two panels side by side, one per cohort (four-survey, six-survey), same axis ranges
%    for direct comparison.
% 3. DATA REQUIRED: Age and elev_percentile_95th from data/processed/master/
%    clean_master_4survey.parquet and clean_master_6survey.parquet (287,064 and 83,382 rows).
% 4. COLOUR/SYMBOL RULES: hexbin density, not scatter (287,064 points would overplot badly);
%    colour scale = point count per hexagon (sequential, single hue per panel); optionally facet
%    or colour-outline by survey year if legible at this density, dropped otherwise in favour of
%    the plain two-panel version.
% 5. ANNOTATIONS: shared x-axis (Age, years) and y-axis (elev_percentile_95th, m) ranges and
%    labels across both panels; a colour-bar legend for point density.
% 6. CAPTION MUST STATE: cohort row/plot counts (already given in Table~\ref{tab:data-funnel}, not
%    repeated numerically here -- caption should describe the pattern in words only, e.g.\ "height
%    increases with age across the observed range," without citing any fitted curve or R^2).
% 7. REMOVE FROM MAIN TEXT: no age/height distribution prose currently exists to remove; this
%    figure adds coverage that was not previously described anywhere in prose.
% 8. DRAFT CAPTION: "Age and ALS-derived top-height coverage for the four-survey (left) and
%    six-survey (right) cohorts, all survey years pooled. Density shown as hexagon bin counts to
%    avoid overplotting; no fitted relationship is shown."
\begin{figure}[!htbp]
\centering
\fbox{\parbox{0.9\textwidth}{\centering\vspace{4cm}FIGURE PLACEHOLDER --- target and age coverage,
fully specified in the comment above source\vspace{4cm}}}
\caption{Age and top-height coverage across cohorts. See source comment for the complete
specification.}
\label{fig:data-target-age}
\end{figure}

\section{Derived terrain, wind, climate and soil variables}
\label{sec:data-environmental}

Every plot in the 71,766-plot environmental table was matched to external terrain, wind, climate
and soil variables, grouped in Table~\ref{tab:data-env-sources} by the same category scheme used
throughout the modelling pipeline. This section summarises each data family's source, resolution
and purpose (Table~\ref{tab:data-env-sources}); the complete per-variable definitions, extraction
methods and coverage statistics are given in \app{Appendix~\ref{app:env-variables}: full
per-variable table}.

% Source material for Appendix~\ref{app:env-variables}: documentation/methodlogy_env_setpick.md
% (verified 2026-08-10) and env_feature_sets_manifest.csv.

% TABLE (final, not a placeholder)
\begin{table}[!htbp]
\centering
\scriptsize
\setlength{\tabcolsep}{3pt}
\renewcommand{\arraystretch}{0.95}
\begin{tabular}{p{1.7cm}p{2.2cm}p{1.8cm}p{2.3cm}p{2.6cm}p{2.9cm}}
\hline
Data family & Source & Resolution & Temporal character & Main purpose & Citation \\
\hline
Forest plots & Forest Research & $\sim$20\,m $\times$ 20\,m polygons & Six survey years, 2002--2023 & Response variable, stand fields & -- \\
Terrain & OS Terrain 50 DTM, Edition 10 & 50\,m grid & Static & Elevation, slope, aspect, exposure & \ct{Ordnance Survey, OS Terrain 50 -- edition/date confirmed for tile NN40 only, see table note} \\
Soil (CEH categorical) & CEH natural capital rasters & 50\,m grid & Static & Pedotope, drainage, texture class & \ct{UKCEH, natural capital rasters} \\
Soil pH & SoilGrids v2.0, ISRIC & 250\,m grid & Static, topsoil (0--5\,cm) & Soil acidity screening & \ct{SoilGrids v2.0, ISRIC} \\
Wind climatology & Global Wind Atlas & National raster grid & Long-term climatology, static & General wind exposure & \ct{Global Wind Atlas, DTU Wind Energy / World Bank ESMAP} \\
Observed wind & MIDAS Open stations & 3 station points & Survey-interval summaries & Temporal wind/storm conditions & \ct{Met Office, MIDAS Open} \\
Climate (long-term) & CHELSA v2.1 BIOCLIM+ & 1\,km grid & Static, 1981--2010 & GDD, precipitation & \ct{CHELSA v2.1} \\
Climate (survey-matched) & HadUK-Grid & 1\,km grid & Time-varying, cohort-averaged & Temperature, ground frost & \ct{Met Office / CEDA, HadUK-Grid} \\
Spatial position & OS Open Roads/Rivers; GPKG polygons & Vector distance & Static & Edge/context effects & \ct{Ordnance Survey, OS Open Roads / OS Open Rivers} \\
\hline
\end{tabular}
\caption{External data sources matched to each plot. Plot density relative to each resolution is
shown in Figure~\ref{fig:data-resolution}, not repeated here. ERA5-Land temperature (8 distinct
values across all 71,766 plots, coarser than HadUK-Grid's 149) was investigated and not retained
in any final feature set. OS Terrain 50 edition/flying date is confirmed for tile NN40 only
(Edition 10, source imagery 2023-05-29, product created 2025-12-25); the other five tiles used
(NN30, NN50, NS39, NS49, NS59) came from the same bulk download but lack individual metadata files.}
\label{tab:data-env-sources}
\end{table}

Global Wind Atlas, TOPEX and MIDAS station records are three distinct, non-interchangeable
representations of wind exposure, even though all three are, in some sense, about wind: Global
Wind Atlas gives a long-term average climatology from a national wind-flow model, not a
measurement at any specific time; TOPEX is a purely geometric calculation from terrain shape
(horizon angles around each plot) describing shelter or exposure, not a wind speed at all; MIDAS
gives real observed wind measurements, but only at three stations several kilometres from
Aberfoyle, summarised over each survey interval. A spatial climatology, a static landscape
property, and an off-site time series answer different questions -- sharing the theme "wind" does
not make them redundant with each other. TOPEX's inclusion is not just convenient because it is
cheap to compute from terrain already in hand: topographic exposure of this kind is an established
predictor of Sitka spruce productivity on Scottish upland sites \citep{worrell1987}, giving this
specific variable a direct forestry justification that the other terrain metrics do not
individually have. Global Wind Atlas, by contrast, was built for wind-energy resource assessment,
not forestry -- its long-term average wind climatology is a reasonable proxy for general spatial
exposure, but was not designed or validated against tree growth or windthrow outcomes, and that
gap between original purpose and this project's use is not resolved by resolution or coverage
alone.

Climate has two different temporal treatments side by side: CHELSA is a single static 1981--2010
climatology applied regardless of a plot's actual survey year (a disclosed temporal mismatch),
while HadUK-Grid's temperature and ground-frost variables are averaged over each cohort's own real
survey years -- closer to survey-year-matched, but still a multi-year mean, not one matched year.
The two answer different questions despite both being climate sources: CHELSA is a global product,
downscaled using terrain to sharpen a coarser interpolation, useful for broad climatic
positioning; HadUK-Grid is built directly from the UK's own observed weather-station network,
closer to a real record of UK conditions than a downscaled global model, which is why it, not
CHELSA, is the one averaged over each cohort's actual survey years.

The circular aspect bearing recorded for terrain is not used directly as a variable: it is
replaced by its two continuous trigonometric components,
$\mathrm{northness}=\cos\theta$ and $\mathrm{eastness}=\sin\theta$, avoiding the discontinuity a
raw 0--360\textdegree{} bearing would otherwise create at north (\app{full derivation and the
frost-hollow flag's exact threshold rule}).

% FIGURE PLACEHOLDER (resolution vs. plot scale) -- fully specified
% 1. PURPOSE: make concrete, in one image, that "50 m" / "250 m" / "1 km" resolution does not mean
%    each plot has an independently measured environmental value -- replacing a paragraph of
%    distinct-value counts with a direct visual comparison against the 400 m^2 plot.
% 2. PANELS: single schematic (not a real map -- a nested-grid diagram at true relative scale).
% 3. DATA REQUIRED: none beyond the numbers already in Table~\ref{tab:data-env-sources} ("plots
%    per grid cell" column) and the plot area fact from \S\ref{sec:data-source} (400 m^2 median).
% 4. COLOUR/SYMBOL RULES: three or four concentric/overlaid squares at true relative scale (50 m,
%    250 m, 1 km, drawn to scale against a 20 m x 20 m plot square at the centre); a small dot
%    cluster inside each ring sized/coloured to represent the "plots per cell" average.
% 5. ANNOTATIONS: each ring labelled with its resolution and data family (terrain/CEH = 50 m,
%    SoilGrids = 250 m, HadUK-Grid/CHELSA = 1 km); plot-count annotation per ring (~4.3, ~53,
%    ~482).
% 6. CAPTION MUST STATE: "a 1 km climate cell covers on average 482 of this study's plots; a 50 m
%    terrain cell covers about 4" -- the two numbers that most concretely convey the mismatch.
% 7. REMOVE FROM MAIN TEXT: the "several hundred plots share one climate cell" sentence from the
%    previous draft's Climate paragraph is now redundant with this figure and should not be
%    restated once the figure is in place.
% 8. DRAFT CAPTION: "Environmental data resolution relative to plot scale. A single 400 m^2 plot
%    sits inside one 50 m terrain cell alongside ~4 neighbouring plots on average, one 250 m soil
%    cell alongside ~53 plots, and one 1 km climate cell alongside ~482 plots. Resolutions are not
%    interchangeable and should not be read as though every plot has an independently measured
%    environmental value."
\begin{figure}[!htbp]
\centering
\fbox{\parbox{0.9\textwidth}{\centering\vspace{4cm}FIGURE PLACEHOLDER --- resolution vs.\ plot
scale schematic, fully specified in the comment above source\vspace{4cm}}}
\caption{Environmental data resolution compared with plot scale. See source comment for the
complete specification.}
\label{fig:data-resolution}
\end{figure}

\section{Strengths, limitations, and what these data can support}
\label{sec:data-strengths-limitations}

Table~\ref{tab:data-strengths-limitations} pairs each notable dataset property with its strength,
its limitation, and the consequence for interpretation, so that no property is presented as purely
positive or purely negative.

% TABLE (final, not a placeholder)
\begin{table}[!htbp]
\centering
\scriptsize
\setlength{\tabcolsep}{3pt}
\renewcommand{\arraystretch}{1.0}
\begin{tabular}{p{2.4cm}p{4.0cm}p{4.0cm}p{4.4cm}}
\hline
Property & Strength & Limitation & Consequence for interpretation \\
\hline
Scale & 71,766 plots (main cohort), 230,359-plot raw layer & -- & Supports stable pooled fits and compartment-level cross-validation with meaningful fold sizes \\
Repeated observations & Up to 6 ALS surveys per plot over 21 years & Only 4--6 observations per plot, at unequal 2--9 year intervals & Limits how finely a within-plot trajectory can be estimated; hides within-interval change \\
Cohort design & Six-survey is a complete, nested subset of four-survey (not partial overlap) & Six-survey covers only 48 compartments vs.\ 296 & Six-survey is a genuine but underpowered sensitivity check, not an independent replication \\
Complete-observation requirement & Guarantees balanced trajectories & May favour plots that stayed identifiable/undisturbed throughout & Sample may not represent a uniformly random cross-section of the forest \\
Response variable & Fine, plot-scale ALS-derived top height & Not identical to biological growth; acquisition metadata limited (\S\ref{sec:data-source}) & Apparent change between surveys may partly reflect measurement, not only growth \\
Species scope & Single, commercially dominant species, clean inventory-based filter & Findings do not generalise beyond Sitka spruce & Appropriately scoped, not a defect -- stated as a boundary, not hidden \\
Spatial hierarchy & Explicit compartment/sub-compartment polygons support real spatial CV units & \texttt{scpt} alone is not unique (\S\ref{sec:data-source}) & Any use of \texttt{scpt} alone elsewhere should be checked \\
Environmental breadth & Terrain, wind, climate, soil all matched to $\geq$99.95\% of plots & Resolution 50\,m--1\,km vs.\ 400\,m\textsuperscript{2} plots; mostly static or long-term (Fig.~\ref{fig:data-resolution}) & Better suited to broad spatial conditions than to explaining a specific between-survey change \\
CanopyCover provenance & Independently informative ALS-derived structural metric & Same-flight provenance as the height target (\S\ref{sec:data-response}) & Open, empirically tested question, not assumed either way \\
Observational design & Rich real-world management data & No experimental manipulation of any variable & Any downstream association is a correlation, not a causal effect \\
\hline
\end{tabular}
\caption{Strengths, limitations, and their consequence for interpretation.}
\label{tab:data-strengths-limitations}
\end{table}

Chapter~\ref{ch:methodology} describes how the modelling design addresses several of these
constraints directly -- for example, through compartment-level spatial-block cross-validation and
an explicit feature-screening pipeline -- but no such mechanism is assumed or explained here.
